\documentclass[10pt,a4paper]{amsart}
\usepackage[margin=19mm]{geometry}
\usepackage{amsmath,amssymb,mathtools}
\usepackage[scr=rsfs,cal=euler]{mathalfa}
\usepackage{xfrac}
\usepackage[unicode,hidelinks,bookmarks=false]{hyperref}
\newcommand{\ie}{i.e.\ }
\newcommand{\cf}{cf.\ }
\newcommand{\resp}{respectively\ }
\newcommand{\Subsection}{\subsection}
\providecommand{\nobreakdash}{\nobreak\hbox{-}}
\setlength{\emergencystretch}{2em}
\multlinegap0pt
\usepackage{bm}
\usepackage{bbm}
\usepackage{dsfont}
\usepackage{xspace}
\usepackage{cancel}
\usepackage{mathtools}
\usepackage{amsthm}
\newtheorem{theorem}{Theorem}
\newtheorem{lemma}[theorem]{Lemma}
\DeclareMathOperator{\asc}{asc}
\DeclareMathOperator{\lev}{lev}
\DeclareMathOperator{\lv}{lv}
\DeclareMathOperator{\pk}{pk}
\DeclareMathOperator{\ret}{ret}
\DeclareMathOperator{\rlmin}{rlmin}
\title[Fibonacci Polyominoes: Refined Enumeration and Dyck-Path Bij]{Fibonacci Polyominoes: Refined Enumeration and Dyck-Path Bijections (Lemma 3.1)}
\author{Jean-Luc Baril, Jos\'e L. Ram\'irez, Samuel Ram\'irez}
\date{Source: arXiv:2609.30097v1 (CC BY 4.0)}
\begin{document}
\maketitle

\noindent\emph{Source and licence.} Fibonacci Polyominoes: Refined Enumeration and Dyck-Path Bijections. Version arXiv:2609.30097v1 (CC BY 4.0), distributed under the Creative Commons Attribution 4.0 International licence, as recorded in the arXiv licence metadata for this exact version and shown on the abstract page https://arxiv.org/abs/2609.30097v1. Full rights notice: licenses/arxiv-2609.30097-CC-BY-4.0.txt.\par\smallskip

\noindent\textit{Editorial scope.} Counted unit: the Lemma printed as Lemma 3.1 of the article (source0.tex lines 1016--1032, source label \texttt{lem:decorated-Dyck}), delivered together with the article's own complete original proof of it (source0.tex lines 1033--1234, one \texttt{proof} environment of 202 lines). No mathematics is written, repaired or re-derived: statement and proof are the article's own text line for line. Required context retained uncounted: none. Every definition, display and label of those lines is retained unchanged. Omitted from this package and not redistributed: 1--1015, 1235--2402. Nothing inside a retained excerpt is omitted or altered: every retained body is delivered exactly as the article's lines are written, so no omission receipt of any kind accompanies this package. Citations: the retained excerpts carry no citation command at all, so no bibliography entry of the article is transcribed and no reference is redistributed. Reference adaptation: none was needed and none was made; every reference inside the delivered unit resolves inside it, so no unresolved reference or citation is produced. Printed slips kept literal: no printed word, symbol, hypothesis or number of the counted statement or of its proof was altered. Layout and class: the article's own class is \texttt{amsart} with packages including hyperref, amsmath, babel, inputenc, amssymb, amsthm, graphics, graphicx, enumerate, ytableau, array, bm, cite, a4wide; the wrapper is the lane's pinned \texttt{amsart} editorial wrapper at 10pt on A4, which restates the theorem-like environments the retained text needs and adds the article's own macro definitions that the retained text needs (\texttt{\textbackslash asc}, \texttt{\textbackslash lev}, \texttt{\textbackslash lv}, \texttt{\textbackslash pk}, \texttt{\textbackslash ret}, \texttt{\textbackslash rlmin}), with extra wrapper packages (bm, bbm, dsfont, xspace, cancel, mathtools). The delivered numbering is the unit's own. Assets: no retained span contains a figure environment, a graphics-inclusion command, a PDF-page inclusion, a table or a TikZ picture; no image file is copied into this package, the package declares no asset dependency and no image file is redistributed with it. Licence: version arXiv:2609.30097v1 of this article is distributed under the Creative Commons Attribution 4.0 International licence, as recorded in the arXiv licence metadata for this exact version and shown on the abstract page https://arxiv.org/abs/2609.30097v1. A full rights notice is delivered at licenses/arxiv-2609.30097-CC-BY-4.0.txt. Characters: every retained excerpt is pure ASCII, so the delivered engine, XeLaTeX with its default Latin Modern text font, reports no missing character.
\begin{lemma}\label{lem:decorated-Dyck}
There is a bijection $\Phi$ between the set $\mathcal{A}_{\ell,r}$ of nonnegative sequences $\mathbf a=(a_1,\ldots, a_\ell)$ with  $\operatorname{exc}(\mathbf a)=r$, and the set $\mathcal{D}_{\ell,r}$ of labeled  Dyck paths of
semilength $\ell$ with total label sum equal to $r$.  Under this bijection, if $\mathbf{a}\in\mathcal{A}_{\ell,r}$ and $(\pi,\lambda)=\Phi(\mathbf a)$, then
\[
\asc(\mathbf{a})=\pk(\pi)-1,
\qquad
\lev(\mathbf{a})=\lv(\pi,\lambda),
\qquad
\rlmin(\mathbf{a})=\ret(\pi).
\]
Moreover, the label of the initial up step of $\pi$ is $\min\{a_1,\ldots,a_\ell\}$.  Consequently, 
\[
\lvert\mathcal{A}_{\ell,r}\rvert
=
C_\ell\binom{r+\ell-1}{\ell-1}.
\]
\end{lemma}

\begin{proof}
  We construct recursively the bijection  $\Phi:\mathcal{A}_{\ell,r}\longrightarrow\mathcal{D}_{\ell,r}$.   For convenience, the empty sequence
corresponds to the empty Dyck path $\epsilon$. Assuming $\ell>0$, let $\mathbf{a}=(a_1,\ldots,a_\ell)\in\mathcal{A}_{\ell,r}$, and set $\rho=\min\{a_1,\ldots,a_\ell\}$. Let $k$ be the rightmost index such that $a_k=\rho$. Define
\[
\mathbf{a}^{L}
=
(a_1-\rho,\ldots,a_{k-1}-\rho) \qquad \text{and} \qquad 
\mathbf{a}^{R}
=
(a_{k+1}-\rho-1,\ldots,a_\ell-\rho-1).
\]
Since $k$ is the rightmost occurrence of the minimum, every entry to its
right is at least $\rho+1$. Thus $\mathbf{a}^{L}$ and
$\mathbf{a}^{R}$ are sequences of nonnegative integers.



  Apply the construction recursively to obtain $\Phi(\mathbf{a}^{L})=(\pi_L,\lambda_L)$ and $\Phi(\mathbf{a}^{R})=(\pi_R,\lambda_R)$. We define
\[
\Phi(\mathbf{a})=(\pi,\lambda),
 \text{ where } 
\pi=U\pi_LD\pi_R \text{ and } 
\]
$\lambda$ is the sequence of nonnegative integers consisting of the juxtaposition of $\rho$, $\lambda_L$ and $\lambda_R$.
 The initial up step of
$\pi$ has label $\rho=\min\{a_1,\ldots,a_\ell\}$.

We next show that the construction is invertible. Let $(\pi,\lambda)$ be
a nonempty labeled Dyck path, and write its first-return decomposition as $\pi=U\pi_LD\pi_R$. Let $\rho$ be the label of the initial up step. Recursively recover
$\mathbf{a}^{L}$ and $\mathbf{a}^{R}$ from the labeled paths
$(\pi_L,\lambda_L)$ and $(\pi_R,\lambda_R)$, respectively, and set
\[
\mathbf{a}
=
\bigl(
\rho+\mathbf{a}^{L},
\ \rho,\
\rho+1+\mathbf{a}^{R}
\bigr),
\]
where addition to a sequence is understood componentwise. Every entry to the
left of the central $\rho$ is at least $\rho$, whereas every entry to its
right is at least $\rho+1$. Hence the central entry is the rightmost
occurrence of the minimum of $\mathbf{a}$. It follows that the two
constructions are inverse to each other.

We now verify the statistics in the statement. We do this recursively,
with all sequence statistics understood to be $0$ on the empty sequence.

First, from the decomposition of $\mathbf a$ one obtains
\[
\operatorname{exc}(\mathbf a)
=
\rho+\operatorname{exc}(\mathbf a^L)+\operatorname{exc}(\mathbf a^R).
\]
Indeed, adding $\rho$ to all entries of $\mathbf a^L$ contributes the extra initial term $\rho$
and leaves all subsequent excesses unchanged. The transition from $\mathbf a^L$ to the central
entry $\rho$ does not contribute to $\operatorname{exc}$, while the transition from
$\rho$ to the first entry on its right, when present, gives exactly the
first term in $\operatorname{exc}(\mathbf{a}^{R})$.  By induction,
$\operatorname{exc}(\mathbf{a}^{L})$ and
$\operatorname{exc}(\mathbf{a}^{R})$ are the sums of the labels of
$\pi_L$ and $\pi_R$, respectively. Since the initial up step has label
$\rho$, we have
\[
|\lambda|
=
\rho+|\lambda_L|+|\lambda_R|
=
\operatorname{exc}(\mathbf{a})
=
r.
\]
Thus $\Phi$ restricts to a bijection between $\mathcal{A}_{\ell,r}$ and
Dyck paths of semilength $\ell$ with total label~$r$.

We next consider ascents. No ascent occurs between the last entry of
$\mathbf{a}^{L}$ and the central entry $\rho$. On the other hand, if
$\mathbf{a}^{R}$ is nonempty, then there is always an ascent from $\rho$
to the first entry on its right. Therefore,
\[
\asc(\mathbf{a})
=
\asc(\mathbf{a}^{L})
+
\asc(\mathbf{a}^{R})
+
\mathbf{1}_{\{\mathbf{a}^{R}\neq\emptyset\}}.
\]
Here $\mathbf{1}_{\{A\}}$ denotes the indicator of the condition $A$, equal to $1$ if $A$ holds and $0$ otherwise. 

For the corresponding first-return decomposition
$\pi=U\pi_LD\pi_R$, the initial $U$ and its matching $D$ form a peak
precisely when $\pi_L$ is empty. Therefore,
\[
\pk(\pi)
=
\pk(\pi_L)
+
\pk(\pi_R)
+
\mathbf{1}_{\{\pi_L=\epsilon\}}.
\]
Using the induction hypothesis on $\pi_L$ and $\pi_R$, we obtain
\[
\pk(\pi_L)
=
\asc(\mathbf a^L)
+
\mathbf{1}_{\{\mathbf a^L\neq\emptyset\}},
\qquad
\pk(\pi_R)
=
\asc(\mathbf a^R)
+
\mathbf{1}_{\{\mathbf a^R\neq\emptyset\}}.
\]
Since $\pi_L=\epsilon$ if and only if $\mathbf a^L=\emptyset$, it follows that
\[
\pk(\pi)
=
\asc(\mathbf a^L)
+
\asc(\mathbf a^R)
+
\mathbf{1}_{\{\mathbf a^R\neq\emptyset\}}
+1
=
\asc(\mathbf a)+1.
\]
Equivalently, $\asc(\mathbf{a})=\pk(\pi)-1$.

We now consider strict right-to-left minima. The central entry $\rho$ is
a strict right-to-left minimum because every entry to its right is
strictly larger than $\rho$. No entry to its left can be a strict
right-to-left minimum, since $\rho$ occurs later in the sequence. The strict
right-to-left minima to the right of $\rho$ are precisely those of
$\mathbf{a}^{R}$. Hence $\rlmin(\mathbf{a})
=
1+\rlmin(\mathbf{a}^{R})$. On the path side, the down step following $\pi_L$ is the first return of
$\pi$, and all subsequent returns are exactly the returns of $\pi_R$.
Thus $
\ret(\pi)
=
1+\ret(\pi_R)$. The induction hypothesis now gives $\rlmin(\mathbf{a})=\ret(\pi)$.

It remains to consider levels. A level can occur inside
$\mathbf{a}^{L}$ or $\mathbf{a}^{R}$. There is never a level between
the central entry $\rho$ and the first entry to its right, since the
latter is at least $\rho+1$. Thus an additional level occurs only
between the last entry on the left and $\rho$. This happens precisely
when $\mathbf{a}^{L}$ is nonempty and its last entry is $0$. Therefore,
\[
\lev(\mathbf{a})
=
\lev(\mathbf{a}^{L})
+
\lev(\mathbf{a}^{R})
+
\begin{cases}
1, & \text{if $\mathbf{a}^{L}$ is nonempty and ends in $0$};\\
0, & \text{otherwise}.
\end{cases}
\]

We claim that $\mathbf{a}^{L}$ is nonempty and ends in $0$ if and only if
$\pi_L$ is nonempty and primitive and its initial up step has label $0$.
Indeed, suppose first that $\mathbf{a}^{L}$ ends in $0$. Since all its
entries are nonnegative, its minimum is $0$, so the initial up step of
$\pi_L$ has label $0$. Moreover, the final $0$ is the only strict
right-to-left minimum of $\mathbf{a}^{L}$. By the result just proved,
$\ret(\pi_L)=1$, and hence $\pi_L$ is primitive. Conversely, the same argument in reverse applies: if $\pi_L$ is primitive
and its initial up step has label $0$, then $\mathbf{a}^{L}$ has minimum
$0$ and exactly one strict right-to-left minimum, which forces its final
entry to be $0$.

It follows that the additional term in the recurrence for
$\lev(\mathbf{a})$ is precisely
$\delta(\pi_L,\lambda_L)$. Using the induction hypothesis gives
\[
\lev(\mathbf{a})
=
\lv(\pi_L,\lambda_L)
+
\lv(\pi_R,\lambda_R)
+
\delta(\pi_L,\lambda_L)
=
\lv(\pi,\lambda).
\]


Finally, there are $C_\ell$ Dyck paths of semilength $\ell$. Once such a
path is fixed, assigning nonnegative labels of total sum $r$ to its
$\ell$ up steps is equivalent to a weak composition of $r$ into $\ell$
parts. There are $\binom{r+\ell-1}{\ell-1}$ such labelings. Consequently,
\[
\lvert\mathcal{A}_{\ell,r}\rvert
=
C_\ell\binom{r+\ell-1}{\ell-1}. \qedhere
\]
\end{proof}

\end{document}
